% Source: https://artofproblemsolving.com/community/c3h2428552_2018_amc_10a_problem_18__nonnegative_integers
% Stable identity: aops-topic-2428552-q01
Let $N=\sum_{i=0}^{7}a_i3^i$ and put $b_i=a_i+1\in\{0,1,2\}$. Then
\[N+\sum_{i=0}^{7}3^i=\sum_{i=0}^{7}b_i3^i,\qquad\text{where}\qquad\sum_{i=0}^{7}3^i=\dfrac{3^8-1}{2}=3280.\]
As $(b_7,\dots,b_0)$ runs over all $3^8=6561$ strings of digits $0,1,2$, the number $\sum b_i3^i$ runs over the base-$3$ representations of $0,1,\dots,6560$, each exactly once, because every integer from $0$ to $3^8-1$ has exactly one $8$-digit base-$3$ representation (with leading zeros). Hence $N$ takes each integer value from $-3280$ to $3280$ exactly once. The nonnegative ones are $0,1,\dots,3280$: there are $\boxed{\textbf{(D) } 3281}$ of them.
